\uuid{6giY}
\titre{Étude de suite}
\auteur{bodin ; hanani ; mzari}
\organisation{exo7}
\licence{CC BY-NC-SA 4.0}
\datecreate{2020-01-20}
\niveau{L1}
\module{Analyse}
\chapitre{Suite}
\sousChapitre{Autre}
\difficulte{3}
\theme{Suites, étude de suite}
\exo7qcmid{lille-1.suites.030}

\contenu{
\begin{question}
\qclassification{121.00}
\qtags{motcle=étude de suite}
On admet que $\forall x \in [0,1[$, $\ln (1+x)\leq x\leq -\ln (1-x)$. Soit $\displaystyle u_n=\sum _{k=1}^n\frac{1}{n+k}$, $n\geq 1$. Quelles sont les bonnes réponses ?
\begin{answers}
    \bad{La suite $(u_n)$ est croissante non majorée.}
    \good{Pour tout $n\geq 1$, $\displaystyle \ln \left(\frac{2n+1}{n+1}\right)\leq u_n\leq \ln (2)$.}
    \good{$(u_n)$ est convergente et $\displaystyle \lim _{n\to +\infty}u_n=\ln (2)$.}
    \bad{$\displaystyle \lim _{n\to +\infty}u_n=+\infty$}
\end{answers}
\begin{explanations}
Avec $\displaystyle x=\frac{1}{n+k}$, on aura :
$$\ln(n+k+1)-\ln (n+k)\leq \frac{1}{n+k}\leq \ln(n+k)-\ln (n+k-1).$$
On somme sur $k$ de $1$ à $n$, on obtient :
$$\ln \left(\frac{2n+1}{n+1}\right)\leq u_n\leq \ln (2n)-\ln (n)=\ln (2).$$
Le théorème d'encadrement implique que $(u_n)$ converge et que sa limite est $\ln (2)$.
\end{explanations}
\end{question}
}
